% !PRINTED 573
\weberpage{573}
%<<CONT
und $\tau_f$ ist congruent mit $\tau$. Demnach erhalten wir die
Congruenz:
\[
[\Pi(t, t_1, t_2, \ldots, t_n)]^p \equiv \Pi(t^p, t_1^p, t_2^p, \ldots, t_n^p)\ (\mathrm{mod}\ \mathfrak p).
\]
Damit ist nach \S.\,150, 4. der Satz bewiesen:

2. \gesperrt{Die Form}
\[
\Pi(t) = (t - \tau)(t - \tau_1) \ldots (t - \tau_{f-1})
\]
\gesperrt{ist nach dem Modul $\mathfrak p$ mit einer ganzen und
homogenen Form $f^{\text{ten}}$ Grades in $R$ der Variablen
$t, t_1, t_2, \ldots, t_n$ congruent.}

Diese ganze rationale Form, deren Coëfficienten bis auf Vielfache der
Primzahl $p$ völlig bestimmt sind, wollen wir mit $P(t)$ bezeichnen.

Dann folgt aus (13)
\weq{14}{\Phi(t) \equiv P(t)\, \Phi_0(t)\ (\mathrm{mod}\ \mathfrak p),}
und da nun $\Phi(t)$ eine ganze Function in $R$ ist, so ergiebt sich
durch Erheben zur Potenz $p$:
\[
\begin{aligned}
\Phi(t^p, u^p) &\equiv P(t^p, u^p)\ \Phi_0(t^p, u^p) \\
[\Phi(t, u)]^p &\equiv [P(t, u)]^p\ [\Phi_0(t, u)]^p
\end{aligned}
\qquad (\mathrm{mod}\ \mathfrak p).
\]
Weil aber die linken Seiten nach dem Modul $\mathfrak p$ congruent
sind, so folgt
\[
[P(t)]^p\ [\Phi_0(t, u)^p - \Phi_0(t^p, u^p)] \equiv 0\ (\mathrm{mod}\ \mathfrak p),
\]
und daraus, da $P(t)$ nicht durch $\mathfrak p$ theilbar ist (weil
der Coëfficient von $t^f$ den Werth 1 hat):
\weq{15}{[\Phi_0(t, u)]^p \equiv \Phi_0(t^p, u^p),}
woraus nach \S.\,150, 4. hervorgeht, dass auch $\Phi_0(t)$ mit einer
ganzen rationalen Form nach dem Modul $\mathfrak p$ congruent ist.

Demnach können wir in (14) auch $\Phi_0(t)$ als ganze Form in $R$
annehmen, und dann muss nach \S.\,140, 3. die Congruenz (14) nicht
nur für den Modul $\mathfrak p$, sondern für den Modul $p$ bestehen.
Daraus erhalten wir den Satz:

3. \gesperrt{Unter den Formen $\Phi(t)$, die für $t = \tau$ in ein
durch $\mathfrak p$ theilbares Functional übergehen, ist $P(t)$ vom
Grade $f$ in Bezug auf $t$ die Form niedrigsten Grades, und wenn
$\Phi(t)$ eine beliebige unter ihnen ist, so lässt sich eine ganze
Function $\Phi_0(t)$ in $R$ so bestimmen, dass}
\weq{15}{\Phi(t) \equiv P(t)\, \Phi_0(t)\ (\mathrm{mod}\ p)} % [sic] number (15) reused for a second equation on this page, as printed
\gesperrt{wird.}
